% Insert into paper.tex as a fifth Extended Data figure, after the ED_Fig4 block.
% The label fig:ed5 is currently unused (paper.tex uses fig:ed1..fig:ed4).
% Caption text: paper_prep/figures/ed5_paralog_allelic/caption.md
% Not applied automatically.

\begin{figure}[h!]
  \includegraphics[width=\linewidth]{fig/ExtendedDataFigures/ED_Fig5_paralog_allelic_scope.pdf}
  \caption{\textbf{Allelic, paralogous and community-level divergence in human
  subtelomeres.}
  \textbf{a}, Cumulative distribution of substitution divergence between
  subtelomeric sequences, weighted by aligned base pairs, for homologous copies of
  the same arm (\emph{allelic}, median 0.15\%), copies on the same chromosome
  (0.83\%) and copies on different chromosomes (0.89\%).
  \textbf{b}, Each class median as fold-change over the allelic median, with the
  range occupied by ancient duplicated gene families (100--1,000-fold) for scale.
  \textbf{c}, For each community, median divergence among member arms divided by
  median divergence to non-member ends; 1.0 indicates no internal cohesion.
  \textbf{d}, Per arm pair, median divergence against the fraction of aligned bases
  at $\geq 99.5\%$ identity, coloured by whether both arms share a community.
  \textbf{e}, For each of 45 chromosome ends, the divergence of its closest partner
  on a different chromosome; dotted line, allelic median; solid line, per-end
  median.}
  \label{fig:ed5}
\end{figure}
